%%%%%%%%%%%%%%%%%%%%%%%%%%%%%%%%%%%%%%%%%%%%%%%%%%%%%%%%%%%%
% 10.8 STOCHASTIC OPTIMIZATION
% The Composition of Advanced Mathematics
%%%%%%%%%%%%%%%%%%%%%%%%%%%%%%%%%%%%%%%%%%%%%%%%%%%%%%%%%%%%

\section{Stochastic Optimization}
\label{sec:stochastic-optimization}

\prerequisite{
Optimization fundamentals, probability, random variables, expectation,
conditional probability, convex optimization, nonlinear programming,
dynamic programming, basic statistics, and Monte Carlo methods.
}

\learningobjectives{
By the end of this section, the reader should be able to:
\begin{itemize}
    \item define stochastic optimization;
    \item distinguish deterministic and stochastic optimization;
    \item identify uncertain parameters and random variables in an
          optimization model;
    \item formulate expected-value objectives;
    \item understand risk-neutral optimization;
    \item recognize stochastic constraints;
    \item distinguish static and multistage decisions;
    \item understand first-stage and recourse decisions;
    \item formulate two-stage stochastic programs;
    \item recognize scenario-based models;
    \item construct deterministic equivalents for finite scenario sets;
    \item understand nonanticipativity constraints;
    \item define the expected value of perfect information;
    \item define the value of the stochastic solution;
    \item distinguish here-and-now and wait-and-see decisions;
    \item understand recourse functions;
    \item recognize relatively complete recourse;
    \item understand sample-average approximation;
    \item recognize Monte Carlo estimation of stochastic objectives;
    \item understand stochastic gradients;
    \item formulate stochastic gradient descent;
    \item understand minibatching;
    \item recognize variance in gradient estimates;
    \item understand Robbins--Monro stochastic approximation conceptually;
    \item recognize convergence conditions for stochastic approximation;
    \item distinguish stochastic programming from robust optimization;
    \item understand chance-constrained optimization;
    \item recognize probabilistic feasibility requirements;
    \item understand scenario approximations of chance constraints;
    \item define risk measures conceptually;
    \item distinguish expectation, variance, Value-at-Risk, and
          Conditional Value-at-Risk;
    \item formulate CVaR optimization;
    \item recognize mean--variance optimization;
    \item understand risk-averse stochastic optimization;
    \item recognize multistage stochastic programs;
    \item understand scenario trees;
    \item formulate nonanticipativity on scenario trees;
    \item recognize stochastic dynamic programming connections;
    \item understand decomposition methods conceptually;
    \item recognize Benders and L-shaped methods;
    \item understand stochastic dual dynamic programming conceptually;
    \item recognize distributionally robust optimization;
    \item understand ambiguity sets conceptually;
    \item distinguish aleatoric uncertainty from model uncertainty
          conceptually;
    \item understand sensitivity to probability models;
    \item connect stochastic optimization with finance, supply chains,
          energy, inventory, machine learning, and control.
\end{itemize}
}

%%%%%%%%%%%%%%%%%%%%%%%%%%%%%%%%%%%%%%%%%%%%%%%%%%%%%%%%%%%%
\subsection{What Is Stochastic Optimization?}
\label{subsec:what-is-stochastic-optimization}
%%%%%%%%%%%%%%%%%%%%%%%%%%%%%%%%%%%%%%%%%%%%%%%%%%%%%%%%%%%%

Stochastic optimization studies optimization problems involving uncertain
quantities represented by random variables.

A generic problem has form

\[
\boxed{
\min_{\mathbf{x}\in X}
\mathbb{E}
\left[
F(\mathbf{x},\boldsymbol{\xi})
\right],
}
\]

where

\[
\mathbf{x}
=
\text{decision variables},
\]

and

\[
\boldsymbol{\xi}
=
\text{random data or uncertainty}.
\]

The decision is chosen while accounting for the probability distribution
of uncertain outcomes.

%%%%%%%%%%%%%%%%%%%%%%%%%%%%%%%%%%%%%%%%%%%%%%%%%%%%%%%%%%%%
\subsection{Deterministic Versus Stochastic Optimization}
\label{subsec:deterministic-versus-stochastic-optimization}
%%%%%%%%%%%%%%%%%%%%%%%%%%%%%%%%%%%%%%%%%%%%%%%%%%%%%%%%%%%%

A deterministic optimization problem assumes all model quantities are
known:

\[
\boxed{
\min_x
f(x).
}
\]

A stochastic optimization problem allows some quantities to be random:

\[
\boxed{
\min_x
\mathbb{E}
[
F(x,\xi)
].
}
\]

Thus stochastic optimization explicitly models uncertainty rather than
replacing it immediately by a fixed estimate.

%%%%%%%%%%%%%%%%%%%%%%%%%%%%%%%%%%%%%%%%%%%%%%%%%%%%%%%%%%%%
\subsection{Sources of Uncertainty}
\label{subsec:sources-of-uncertainty}
%%%%%%%%%%%%%%%%%%%%%%%%%%%%%%%%%%%%%%%%%%%%%%%%%%%%%%%%%%%%

Uncertainty can enter through:

\[
\boxed{
\text{demand},
}
\]

\[
\boxed{
\text{prices},
}
\]

\[
\boxed{
\text{travel times},
}
\]

\[
\boxed{
\text{equipment failures},
}
\]

\[
\boxed{
\text{weather},
}
\]

\[
\boxed{
\text{returns},
}
\]

or

\[
\boxed{
\text{measurement noise}.
}
\]

The probability model attempts to describe these uncertain quantities.

%%%%%%%%%%%%%%%%%%%%%%%%%%%%%%%%%%%%%%%%%%%%%%%%%%%%%%%%%%%%
\subsection{Risk-Neutral Objective}
\label{subsec:risk-neutral-objective}
%%%%%%%%%%%%%%%%%%%%%%%%%%%%%%%%%%%%%%%%%%%%%%%%%%%%%%%%%%%%

A risk-neutral decision maker minimizes expected cost:

\[
\boxed{
\min_x
\mathbb{E}
[
C(x,\xi)
].
}
\]

Equivalently, for reward maximization,

\[
\boxed{
\max_x
\mathbb{E}
[
R(x,\xi)
].
}
\]

Only the expected value enters the objective directly.

%%%%%%%%%%%%%%%%%%%%%%%%%%%%%%%%%%%%%%%%%%%%%%%%%%%%%%%%%%%%
\subsection{Expected Objective with Discrete Scenarios}
\label{subsec:expected-objective-discrete-scenarios}
%%%%%%%%%%%%%%%%%%%%%%%%%%%%%%%%%%%%%%%%%%%%%%%%%%%%%%%%%%%%

Suppose uncertainty has scenarios

\[
\omega=1,\ldots,N
\]

with probabilities

\[
p_\omega,
\qquad
p_\omega\geq0,
\]

and

\[
\sum_{\omega=1}^{N}
p_\omega=1.
\]

Then

\[
\boxed{
\mathbb{E}
[
F(x,\xi)
]
=
\sum_{\omega=1}^{N}
p_\omega
F(x,\xi^\omega).
}
\]

%%%%%%%%%%%%%%%%%%%%%%%%%%%%%%%%%%%%%%%%%%%%%%%%%%%%%%%%%%%%
\subsection{Worked Example: Expected Cost}
\label{subsec:worked-expected-cost}
%%%%%%%%%%%%%%%%%%%%%%%%%%%%%%%%%%%%%%%%%%%%%%%%%%%%%%%%%%%%

\begin{workedexamplebox}{Choosing Between Two Plans}

Suppose plan \(A\) has costs

\[
10
\]

with probability

\[
0.8,
\]

and

\[
30
\]

with probability

\[
0.2.
\]

Its expected cost is

\[
0.8(10)+0.2(30)
=
14.
\]

Plan \(B\) always costs

\[
16.
\]

Thus,

\[
\mathbb{E}[C_A]=14,
\]

and

\[
\mathbb{E}[C_B]=16.
\]

A risk-neutral decision maker chooses

\begin{answerbox}
\[
\boxed{
\text{Plan }A.
}
\]
\end{answerbox}

\end{workedexamplebox}

%%%%%%%%%%%%%%%%%%%%%%%%%%%%%%%%%%%%%%%%%%%%%%%%%%%%%%%%%%%%
\subsection{Expectation Does Not Describe All Risk}
\label{subsec:expectation-not-all-risk}
%%%%%%%%%%%%%%%%%%%%%%%%%%%%%%%%%%%%%%%%%%%%%%%%%%%%%%%%%%%%

Two decisions can have the same expected cost but very different
distributions.

For example:

\[
C_A
=
10
\]

with certainty,

while

\[
C_B
=
\begin{cases}
0,
&
\text{with probability }\frac12,
\\
20,
&
\text{with probability }\frac12.
\end{cases}
\]

Both have expectation

\[
\boxed{
10.
}
\]

Yet decision \(B\) is more variable.

This motivates risk-averse optimization.

%%%%%%%%%%%%%%%%%%%%%%%%%%%%%%%%%%%%%%%%%%%%%%%%%%%%%%%%%%%%
\subsection{Here-and-Now Decisions}
\label{subsec:here-and-now-decisions}
%%%%%%%%%%%%%%%%%%%%%%%%%%%%%%%%%%%%%%%%%%%%%%%%%%%%%%%%%%%%

A here-and-now decision must be chosen before uncertainty is observed.

For example:

\[
\boxed{
\text{capacity installed before demand is known}.
}
\]

Such decisions cannot depend on the realized scenario.

They are often called first-stage decisions.

%%%%%%%%%%%%%%%%%%%%%%%%%%%%%%%%%%%%%%%%%%%%%%%%%%%%%%%%%%%%
\subsection{Recourse Decisions}
\label{subsec:recourse-decisions}
%%%%%%%%%%%%%%%%%%%%%%%%%%%%%%%%%%%%%%%%%%%%%%%%%%%%%%%%%%%%

After uncertainty is observed, additional corrective decisions may be
allowed.

These are recourse decisions.

Examples include:

\[
\boxed{
\text{overtime},
}
\]

\[
\boxed{
\text{emergency purchasing},
}
\]

\[
\boxed{
\text{rerouting},
}
\]

or

\[
\boxed{
\text{inventory adjustment}.
}
\]

%%%%%%%%%%%%%%%%%%%%%%%%%%%%%%%%%%%%%%%%%%%%%%%%%%%%%%%%%%%%
\subsection{Two-Stage Stochastic Programming}
\label{subsec:two-stage-stochastic-programming}
%%%%%%%%%%%%%%%%%%%%%%%%%%%%%%%%%%%%%%%%%%%%%%%%%%%%%%%%%%%%

A two-stage stochastic program has:

\[
\boxed{
\text{first-stage decisions before uncertainty},
}
\]

followed by

\[
\boxed{
\text{second-stage decisions after uncertainty}.
}
\]

A standard form is

\[
\boxed{
\min_x
\left[
c^Tx
+
\mathbb{E}_{\xi}
[
Q(x,\xi)
]
\right],
}
\]

subject to first-stage constraints.

%%%%%%%%%%%%%%%%%%%%%%%%%%%%%%%%%%%%%%%%%%%%%%%%%%%%%%%%%%%%
\subsection{Recourse Function}
\label{subsec:recourse-function}
%%%%%%%%%%%%%%%%%%%%%%%%%%%%%%%%%%%%%%%%%%%%%%%%%%%%%%%%%%%%

The second-stage recourse value can be defined as

\[
\boxed{
Q(x,\xi)
=
\min_y
\left\{
q(\xi)^Ty:
W(\xi)y
=
h(\xi)-T(\xi)x,
\;
y\geq0
\right\}.
}
\]

The first-stage problem becomes

\[
\boxed{
\min_x
c^Tx
+
\mathbb{E}
[
Q(x,\xi)
].
}
\]

%%%%%%%%%%%%%%%%%%%%%%%%%%%%%%%%%%%%%%%%%%%%%%%%%%%%%%%%%%%%
\subsection{Meaning of Recourse}
\label{subsec:meaning-of-recourse}
%%%%%%%%%%%%%%%%%%%%%%%%%%%%%%%%%%%%%%%%%%%%%%%%%%%%%%%%%%%%

The variable \(x\) prepares for uncertain future conditions.

After scenario \(\xi\) occurs, the recourse decision \(y\) adjusts the
system.

Thus stochastic programming balances:

\[
\boxed{
\text{cost of preparation}
}
\]

against

\[
\boxed{
\text{expected cost of future correction}.
}
\]

%%%%%%%%%%%%%%%%%%%%%%%%%%%%%%%%%%%%%%%%%%%%%%%%%%%%%%%%%%%%
\subsection{Worked Example: Simple Recourse}
\label{subsec:worked-simple-recourse}
%%%%%%%%%%%%%%%%%%%%%%%%%%%%%%%%%%%%%%%%%%%%%%%%%%%%%%%%%%%%

\begin{workedexamplebox}{Capacity and Emergency Supply}

Suppose a company chooses capacity \(x\) before demand \(D\) is known.

Installed capacity costs

\[
2x.
\]

If demand exceeds capacity, emergency supply costs

\[
5
\]

per unit.

The recourse quantity is

\[
y(D)
=
\max
\{
D-x,0
\}.
\]

Therefore the stochastic objective is

\[
\boxed{
\min_{x\geq0}
\left[
2x
+
5
\mathbb{E}
\left[
(D-x)_+
\right]
\right],
}
\]

where

\[
(u)_+
=
\max\{u,0\}.
\]

This model trades fixed capacity against expected shortage cost.

\end{workedexamplebox}

%%%%%%%%%%%%%%%%%%%%%%%%%%%%%%%%%%%%%%%%%%%%%%%%%%%%%%%%%%%%
\subsection{Finite-Scenario Deterministic Equivalent}
\label{subsec:finite-scenario-deterministic-equivalent}
%%%%%%%%%%%%%%%%%%%%%%%%%%%%%%%%%%%%%%%%%%%%%%%%%%%%%%%%%%%%

Suppose there are scenarios

\[
\omega=1,\ldots,N.
\]

For each scenario, introduce second-stage variables

\[
y^\omega.
\]

Then

\[
\boxed{
\begin{aligned}
\min_{x,\{y^\omega\}}
\quad&
c^Tx
+
\sum_{\omega=1}^{N}
p_\omega
q_\omega^T
y^\omega
\\
\text{subject to}
\quad&
Ax=b,
\\
&
T_\omega x
+
W_\omega y^\omega
=
h_\omega,
\qquad
\omega=1,\ldots,N,
\\
&
x\geq0,
\qquad
y^\omega\geq0.
\end{aligned}
}
\]

This is called the deterministic equivalent.

%%%%%%%%%%%%%%%%%%%%%%%%%%%%%%%%%%%%%%%%%%%%%%%%%%%%%%%%%%%%
\subsection{Scenario Expansion}
\label{subsec:scenario-expansion}
%%%%%%%%%%%%%%%%%%%%%%%%%%%%%%%%%%%%%%%%%%%%%%%%%%%%%%%%%%%%

Each scenario receives its own recourse variables.

Thus the number of variables can increase approximately in proportion to
the number of scenarios.

Large scenario sets can therefore produce very large optimization models.

This motivates decomposition and sampling methods.

%%%%%%%%%%%%%%%%%%%%%%%%%%%%%%%%%%%%%%%%%%%%%%%%%%%%%%%%%%%%
\subsection{Nonanticipativity}
\label{subsec:nonanticipativity}
%%%%%%%%%%%%%%%%%%%%%%%%%%%%%%%%%%%%%%%%%%%%%%%%%%%%%%%%%%%%

A decision cannot depend on information that has not yet been observed.

This is the principle of nonanticipativity.

If two scenarios are indistinguishable at stage \(t\), then the decisions
made at that stage must be identical.

Symbolically,

\[
\boxed{
x_t^\omega
=
x_t^{\omega'}
}
\]

whenever scenarios \(\omega\) and \(\omega'\) share the same history up to
stage \(t\).

%%%%%%%%%%%%%%%%%%%%%%%%%%%%%%%%%%%%%%%%%%%%%%%%%%%%%%%%%%%%
\subsection{Scenario Trees}
\label{subsec:scenario-trees}
%%%%%%%%%%%%%%%%%%%%%%%%%%%%%%%%%%%%%%%%%%%%%%%%%%%%%%%%%%%%

A multistage stochastic process can be represented by a scenario tree.

Each node represents the information available at a stage.

Branches represent possible future realizations.

A decision at a node depends only on the history leading to that node.

Thus scenario trees encode nonanticipativity structurally.

%%%%%%%%%%%%%%%%%%%%%%%%%%%%%%%%%%%%%%%%%%%%%%%%%%%%%%%%%%%%
\subsection{Multistage Stochastic Programming}
\label{subsec:multistage-stochastic-programming}
%%%%%%%%%%%%%%%%%%%%%%%%%%%%%%%%%%%%%%%%%%%%%%%%%%%%%%%%%%%%

A multistage problem alternates between decisions and observations:

\[
\boxed{
x_0
\rightarrow
\xi_1
\rightarrow
x_1
\rightarrow
\xi_2
\rightarrow
x_2
\rightarrow
\cdots.
}
\]

Each decision \(x_t\) may depend on uncertainty observed through stage
\(t\), but not on future realizations.

%%%%%%%%%%%%%%%%%%%%%%%%%%%%%%%%%%%%%%%%%%%%%%%%%%%%%%%%%%%%
\subsection{Nested Expectation Form}
\label{subsec:nested-expectation-stochastic}
%%%%%%%%%%%%%%%%%%%%%%%%%%%%%%%%%%%%%%%%%%%%%%%%%%%%%%%%%%%%

A multistage objective may be written schematically as

\[
\boxed{
\min
\mathbb{E}
\left[
c_0(x_0)
+
c_1(x_1,\xi_1)
+
\cdots
+
c_T(x_T,\xi_T)
\right].
}
\]

The decision rule at each stage must respect the information available at
that stage.

%%%%%%%%%%%%%%%%%%%%%%%%%%%%%%%%%%%%%%%%%%%%%%%%%%%%%%%%%%%%
\subsection{Stochastic Dynamic Programming Connection}
\label{subsec:stochastic-dp-connection}
%%%%%%%%%%%%%%%%%%%%%%%%%%%%%%%%%%%%%%%%%%%%%%%%%%%%%%%%%%%%

Multistage stochastic optimization can often be written through Bellman
recursions.

For state \(s_t\),

\[
\boxed{
V_t(s_t)
=
\min_{a_t}
\left[
c_t(s_t,a_t)
+
\mathbb{E}
\left[
V_{t+1}(S_{t+1})
\mid
s_t,a_t
\right]
\right].
}
\]

Thus stochastic programming and stochastic dynamic programming are closely
related.

%%%%%%%%%%%%%%%%%%%%%%%%%%%%%%%%%%%%%%%%%%%%%%%%%%%%%%%%%%%%
\subsection{Wait-and-See Solution}
\label{subsec:wait-and-see-solution}
%%%%%%%%%%%%%%%%%%%%%%%%%%%%%%%%%%%%%%%%%%%%%%%%%%%%%%%%%%%%

The wait-and-see problem assumes future uncertainty is known before the
decision is made.

For each scenario \(\xi\), solve

\[
\boxed{
\min_x
F(x,\xi).
}
\]

Then average the optimal scenario values:

\[
\boxed{
WS
=
\mathbb{E}
\left[
\min_x
F(x,\xi)
\right].
}
\]

This represents perfect information.

%%%%%%%%%%%%%%%%%%%%%%%%%%%%%%%%%%%%%%%%%%%%%%%%%%%%%%%%%%%%
\subsection{Stochastic Solution}
\label{subsec:stochastic-solution}
%%%%%%%%%%%%%%%%%%%%%%%%%%%%%%%%%%%%%%%%%%%%%%%%%%%%%%%%%%%%

The true stochastic solution solves

\[
\boxed{
RP
=
\min_x
\mathbb{E}
[
F(x,\xi)
],
}
\]

where \(RP\) is often called the recourse problem value in standard
two-stage terminology.

The decision must be chosen before the uncertain realization is known.

%%%%%%%%%%%%%%%%%%%%%%%%%%%%%%%%%%%%%%%%%%%%%%%%%%%%%%%%%%%%
\subsection{Expected Value Problem}
\label{subsec:expected-value-problem}
%%%%%%%%%%%%%%%%%%%%%%%%%%%%%%%%%%%%%%%%%%%%%%%%%%%%%%%%%%%%

A common approximation replaces the random quantity by its expectation:

\[
\boxed{
\bar{\xi}
=
\mathbb{E}[\xi].
}
\]

Then solve

\[
\boxed{
\min_x
F(x,\bar{\xi}).
}
\]

Let the resulting decision be

\[
\bar{x}.
\]

This is the expected-value solution.

%%%%%%%%%%%%%%%%%%%%%%%%%%%%%%%%%%%%%%%%%%%%%%%%%%%%%%%%%%%%
\subsection{Expected Result of the Expected-Value Solution}
\label{subsec:eev}
%%%%%%%%%%%%%%%%%%%%%%%%%%%%%%%%%%%%%%%%%%%%%%%%%%%%%%%%%%%%

Evaluate the expected-value decision under the actual stochastic model:

\[
\boxed{
EEV
=
\mathbb{E}
[
F(\bar{x},\xi)
].
}
\]

This measures how the deterministic expected-value solution performs under
uncertainty.

%%%%%%%%%%%%%%%%%%%%%%%%%%%%%%%%%%%%%%%%%%%%%%%%%%%%%%%%%%%%
\subsection{Value of the Stochastic Solution}
\label{subsec:value-of-stochastic-solution}
%%%%%%%%%%%%%%%%%%%%%%%%%%%%%%%%%%%%%%%%%%%%%%%%%%%%%%%%%%%%

For a minimization problem,

\[
\boxed{
VSS
=
EEV-RP.
}
\]

Since the stochastic solution is optimized over the full uncertainty
model,

\[
\boxed{
VSS\geq0
}
\]

under the standard definitions.

A large VSS means explicit stochastic modeling provides substantial value.

%%%%%%%%%%%%%%%%%%%%%%%%%%%%%%%%%%%%%%%%%%%%%%%%%%%%%%%%%%%%
\subsection{Expected Value of Perfect Information}
\label{subsec:expected-value-perfect-information}
%%%%%%%%%%%%%%%%%%%%%%%%%%%%%%%%%%%%%%%%%%%%%%%%%%%%%%%%%%%%

For minimization,

\[
\boxed{
EVPI
=
RP-WS.
}
\]

Because perfect information cannot make the optimum worse,

\[
\boxed{
WS\leq RP.
}
\]

Thus

\[
\boxed{
EVPI\geq0.
}
\]

EVPI is the maximum amount one should be willing to pay for perfect
information under the model.

%%%%%%%%%%%%%%%%%%%%%%%%%%%%%%%%%%%%%%%%%%%%%%%%%%%%%%%%%%%%
\subsection{Worked Example: EVPI}
\label{subsec:worked-evpi}
%%%%%%%%%%%%%%%%%%%%%%%%%%%%%%%%%%%%%%%%%%%%%%%%%%%%%%%%%%%%

\begin{workedexamplebox}{Value of Perfect Information}

Suppose the optimal stochastic expected cost is

\[
RP=120.
\]

If future uncertainty were known perfectly before deciding, suppose the
expected wait-and-see cost would be

\[
WS=100.
\]

Then

\[
EVPI
=
120-100.
\]

Therefore,

\begin{answerbox}
\[
\boxed{
EVPI=20.
}
\]
\end{answerbox}

Under the model, perfect information is worth at most \(20\) units of
expected cost reduction.

\end{workedexamplebox}

%%%%%%%%%%%%%%%%%%%%%%%%%%%%%%%%%%%%%%%%%%%%%%%%%%%%%%%%%%%%
\subsection{Relatively Complete Recourse}
\label{subsec:relatively-complete-recourse}
%%%%%%%%%%%%%%%%%%%%%%%%%%%%%%%%%%%%%%%%%%%%%%%%%%%%%%%%%%%%

A two-stage model has relatively complete recourse if every first-stage
decision satisfying the first-stage constraints admits a feasible
second-stage response for almost every scenario.

Conceptually,

\[
\boxed{
\text{every admissible first-stage decision can be repaired later}.
}
\]

Without this property, some first-stage decisions may lead to infeasible
future scenarios.

%%%%%%%%%%%%%%%%%%%%%%%%%%%%%%%%%%%%%%%%%%%%%%%%%%%%%%%%%%%%
\subsection{Complete Recourse}
\label{subsec:complete-recourse}
%%%%%%%%%%%%%%%%%%%%%%%%%%%%%%%%%%%%%%%%%%%%%%%%%%%%%%%%%%%%

Complete recourse is a stronger structural condition on the second-stage
system.

Roughly, the recourse matrix permits a feasible second-stage correction
for every required right-hand-side vector.

Precise definitions depend on the recourse model.

%%%%%%%%%%%%%%%%%%%%%%%%%%%%%%%%%%%%%%%%%%%%%%%%%%%%%%%%%%%%
\subsection{Scenario Generation}
\label{subsec:scenario-generation}
%%%%%%%%%%%%%%%%%%%%%%%%%%%%%%%%%%%%%%%%%%%%%%%%%%%%%%%%%%%%

When uncertainty is continuous or extremely large, one often approximates
it with a finite collection of scenarios.

Scenarios may be generated from:

\[
\boxed{
\text{historical data},
}
\]

\[
\boxed{
\text{probability models},
}
\]

\[
\boxed{
\text{Monte Carlo sampling},
}
\]

or

\[
\boxed{
\text{expert-defined stress cases}.
}
\]

%%%%%%%%%%%%%%%%%%%%%%%%%%%%%%%%%%%%%%%%%%%%%%%%%%%%%%%%%%%%
\subsection{Scenario Reduction}
\label{subsec:scenario-reduction}
%%%%%%%%%%%%%%%%%%%%%%%%%%%%%%%%%%%%%%%%%%%%%%%%%%%%%%%%%%%%

A large scenario set may be compressed to a smaller representative set.

Scenario reduction attempts to preserve the important features of the
uncertainty distribution while reducing computational burden.

This creates a tradeoff between:

\[
\boxed{
\text{statistical fidelity}
}
\]

and

\[
\boxed{
\text{computational cost}.
}
\]

%%%%%%%%%%%%%%%%%%%%%%%%%%%%%%%%%%%%%%%%%%%%%%%%%%%%%%%%%%%%
\subsection{Sample-Average Approximation}
\label{subsec:sample-average-approximation}
%%%%%%%%%%%%%%%%%%%%%%%%%%%%%%%%%%%%%%%%%%%%%%%%%%%%%%%%%%%%

Suppose

\[
\xi^{(1)},
\ldots,
\xi^{(N)}
\]

are independent samples from the uncertainty distribution.

The true problem

\[
\boxed{
\min_x
\mathbb{E}
[
F(x,\xi)
]
}
\]

is approximated by

\[
\boxed{
\min_x
\frac1N
\sum_{i=1}^{N}
F
(
x,\xi^{(i)}
).
}
\]

This is called sample-average approximation, abbreviated SAA.

%%%%%%%%%%%%%%%%%%%%%%%%%%%%%%%%%%%%%%%%%%%%%%%%%%%%%%%%%%%%
\subsection{Law of Large Numbers Connection}
\label{subsec:saa-law-large-numbers}
%%%%%%%%%%%%%%%%%%%%%%%%%%%%%%%%%%%%%%%%%%%%%%%%%%%%%%%%%%%%

For a fixed \(x\), under suitable conditions,

\[
\boxed{
\frac1N
\sum_{i=1}^{N}
F
(
x,\xi^{(i)}
)
\to
\mathbb{E}
[
F(x,\xi)
]
}
\]

as

\[
N\to\infty.
\]

This follows from laws of large numbers.

Optimization adds additional questions about convergence of minimizers and
optimal values.

%%%%%%%%%%%%%%%%%%%%%%%%%%%%%%%%%%%%%%%%%%%%%%%%%%%%%%%%%%%%
\subsection{Worked Example: Sample Average}
\label{subsec:worked-sample-average}
%%%%%%%%%%%%%%%%%%%%%%%%%%%%%%%%%%%%%%%%%%%%%%%%%%%%%%%%%%%%

\begin{workedexamplebox}{Monte Carlo Objective Estimate}

Suppose a decision \(x\) produces sampled costs

\[
12,\quad
15,\quad
9,\quad
14,\quad
10.
\]

The sample-average estimate is

\[
\frac{
12+15+9+14+10
}{
5
}
=
12.
\]

Therefore,

\begin{answerbox}
\[
\boxed{
\widehat{\mathbb{E}[C(x,\xi)]}
=
12.
}
\]
\end{answerbox}

\end{workedexamplebox}

%%%%%%%%%%%%%%%%%%%%%%%%%%%%%%%%%%%%%%%%%%%%%%%%%%%%%%%%%%%%
\subsection{Out-of-Sample Evaluation}
\label{subsec:out-of-sample-evaluation-stochastic}
%%%%%%%%%%%%%%%%%%%%%%%%%%%%%%%%%%%%%%%%%%%%%%%%%%%%%%%%%%%%

A solution optimized on one sample may fit that sample too closely.

Therefore it should be evaluated on independent scenarios not used during
optimization.

This is called out-of-sample evaluation.

It estimates how the decision performs under new realizations of
uncertainty.

%%%%%%%%%%%%%%%%%%%%%%%%%%%%%%%%%%%%%%%%%%%%%%%%%%%%%%%%%%%%
\subsection{Training and Validation Analogy}
\label{subsec:training-validation-stochastic}
%%%%%%%%%%%%%%%%%%%%%%%%%%%%%%%%%%%%%%%%%%%%%%%%%%%%%%%%%%%%

Sample-average optimization resembles statistical learning:

\[
\boxed{
\text{optimization sample}
\leftrightarrow
\text{training data},
}
\]

and

\[
\boxed{
\text{evaluation sample}
\leftrightarrow
\text{validation or test data}.
}
\]

This analogy is especially important in data-driven optimization.

%%%%%%%%%%%%%%%%%%%%%%%%%%%%%%%%%%%%%%%%%%%%%%%%%%%%%%%%%%%%
\subsection{Stochastic Approximation}
\label{subsec:stochastic-approximation}
%%%%%%%%%%%%%%%%%%%%%%%%%%%%%%%%%%%%%%%%%%%%%%%%%%%%%%%%%%%%

Suppose the objective is

\[
\boxed{
f(x)
=
\mathbb{E}
[
F(x,\xi)
],
}
\]

but evaluating the expectation exactly is expensive.

If an unbiased gradient estimate is available,

\[
\boxed{
\mathbb{E}
[
G(x,\xi)
]
=
\nabla f(x),
}
\]

one can update using random samples.

%%%%%%%%%%%%%%%%%%%%%%%%%%%%%%%%%%%%%%%%%%%%%%%%%%%%%%%%%%%%
\subsection{Stochastic Gradient Descent}
\label{subsec:stochastic-gradient-descent}
%%%%%%%%%%%%%%%%%%%%%%%%%%%%%%%%%%%%%%%%%%%%%%%%%%%%%%%%%%%%

Stochastic gradient descent, abbreviated SGD, uses

\[
\boxed{
x_{k+1}
=
x_k
-
\alpha_k
G(x_k,\xi_k),
}
\]

where

\[
\xi_k
\]

is a random sample and

\[
G(x_k,\xi_k)
\]

is a stochastic gradient estimate.

%%%%%%%%%%%%%%%%%%%%%%%%%%%%%%%%%%%%%%%%%%%%%%%%%%%%%%%%%%%%
\subsection{Unbiased Stochastic Gradient}
\label{subsec:unbiased-stochastic-gradient}
%%%%%%%%%%%%%%%%%%%%%%%%%%%%%%%%%%%%%%%%%%%%%%%%%%%%%%%%%%%%

An ideal stochastic gradient satisfies

\[
\boxed{
\mathbb{E}
[
G(x,\xi)
]
=
\nabla f(x).
}
\]

Individual gradient estimates can be noisy.

Their average direction, however, matches the true gradient.

%%%%%%%%%%%%%%%%%%%%%%%%%%%%%%%%%%%%%%%%%%%%%%%%%%%%%%%%%%%%
\subsection{Worked Example: Stochastic Gradient}
\label{subsec:worked-stochastic-gradient}
%%%%%%%%%%%%%%%%%%%%%%%%%%%%%%%%%%%%%%%%%%%%%%%%%%%%%%%%%%%%

\begin{workedexamplebox}{One SGD Step}

Suppose

\[
x_k=4,
\]

the sampled stochastic gradient is

\[
G(x_k,\xi_k)=6,
\]

and

\[
\alpha_k=0.1.
\]

Then

\[
x_{k+1}
=
4-0.1(6).
\]

Thus,

\begin{answerbox}
\[
\boxed{
x_{k+1}=3.4.
}
\]
\end{answerbox}

\end{workedexamplebox}

%%%%%%%%%%%%%%%%%%%%%%%%%%%%%%%%%%%%%%%%%%%%%%%%%%%%%%%%%%%%
\subsection{Minibatches}
\label{subsec:minibatches-stochastic-optimization}
%%%%%%%%%%%%%%%%%%%%%%%%%%%%%%%%%%%%%%%%%%%%%%%%%%%%%%%%%%%%

Instead of using one sample, use a minibatch

\[
B_k.
\]

The gradient estimate is

\[
\boxed{
G_k
=
\frac1{|B_k|}
\sum_{i\in B_k}
\nabla
F
(
x_k,\xi_i
).
}
\]

Larger batches generally reduce gradient variance but require more work per
iteration.

%%%%%%%%%%%%%%%%%%%%%%%%%%%%%%%%%%%%%%%%%%%%%%%%%%%%%%%%%%%%
\subsection{Variance of Gradient Estimates}
\label{subsec:gradient-estimate-variance}
%%%%%%%%%%%%%%%%%%%%%%%%%%%%%%%%%%%%%%%%%%%%%%%%%%%%%%%%%%%%

Suppose

\[
\mathbb{E}
[
G(x,\xi)
]
=
\nabla f(x).
\]

The noise is

\[
\boxed{
G(x,\xi)
-
\nabla f(x).
}
\]

Its variance affects convergence.

High-variance estimates can cause highly irregular iterates.

%%%%%%%%%%%%%%%%%%%%%%%%%%%%%%%%%%%%%%%%%%%%%%%%%%%%%%%%%%%%
\subsection{Step-Size Sequences}
\label{subsec:stochastic-step-size-sequences}
%%%%%%%%%%%%%%%%%%%%%%%%%%%%%%%%%%%%%%%%%%%%%%%%%%%%%%%%%%%%

Classical stochastic approximation often uses decreasing step sizes
satisfying conditions such as

\[
\boxed{
\sum_{k=1}^{\infty}
\alpha_k
=
\infty,
}
\]

and

\[
\boxed{
\sum_{k=1}^{\infty}
\alpha_k^2
<
\infty.
}
\]

A common example is approximately

\[
\boxed{
\alpha_k
=
\frac{c}{k}.
}
\]

Exact convergence requirements depend on the problem assumptions.

%%%%%%%%%%%%%%%%%%%%%%%%%%%%%%%%%%%%%%%%%%%%%%%%%%%%%%%%%%%%
\subsection{Robbins--Monro Stochastic Approximation}
\label{subsec:robbins-monro}
%%%%%%%%%%%%%%%%%%%%%%%%%%%%%%%%%%%%%%%%%%%%%%%%%%%%%%%%%%%%

The Robbins--Monro method was introduced for solving equations involving
expectations when only noisy observations are available.

A generic stochastic approximation update is

\[
\boxed{
x_{k+1}
=
x_k
-
\alpha_k
G_k.
}
\]

Modern stochastic gradient methods are closely related to this framework.

%%%%%%%%%%%%%%%%%%%%%%%%%%%%%%%%%%%%%%%%%%%%%%%%%%%%%%%%%%%%
\subsection{Constant Step Sizes}
\label{subsec:constant-step-size-sgd}
%%%%%%%%%%%%%%%%%%%%%%%%%%%%%%%%%%%%%%%%%%%%%%%%%%%%%%%%%%%%

A constant step size can produce rapid early progress.

However, persistent gradient noise may prevent exact convergence to a
single optimum.

Instead, iterates may fluctuate in a neighborhood.

Decreasing step sizes can reduce these fluctuations asymptotically.

%%%%%%%%%%%%%%%%%%%%%%%%%%%%%%%%%%%%%%%%%%%%%%%%%%%%%%%%%%%%
\subsection{Iterate Averaging}
\label{subsec:iterate-averaging}
%%%%%%%%%%%%%%%%%%%%%%%%%%%%%%%%%%%%%%%%%%%%%%%%%%%%%%%%%%%%

An averaged iterate may be defined by

\[
\boxed{
\bar{x}_N
=
\frac1N
\sum_{k=1}^{N}
x_k.
}
\]

Averaging can improve stability and theoretical convergence properties in
stochastic approximation.

%%%%%%%%%%%%%%%%%%%%%%%%%%%%%%%%%%%%%%%%%%%%%%%%%%%%%%%%%%%%
\subsection{Variance Reduction Preview}
\label{subsec:stochastic-variance-reduction}
%%%%%%%%%%%%%%%%%%%%%%%%%%%%%%%%%%%%%%%%%%%%%%%%%%%%%%%%%%%%

For finite-sum objectives,

\[
\boxed{
f(x)
=
\frac1N
\sum_{i=1}^{N}
f_i(x),
}
\]

variance-reduced stochastic methods attempt to retain the low per-step cost
of SGD while reducing stochastic gradient noise.

Examples include methods based on periodically refreshed full-gradient
information or stored component gradients.

%%%%%%%%%%%%%%%%%%%%%%%%%%%%%%%%%%%%%%%%%%%%%%%%%%%%%%%%%%%%
\subsection{Chance-Constrained Optimization}
\label{subsec:chance-constrained-optimization}
%%%%%%%%%%%%%%%%%%%%%%%%%%%%%%%%%%%%%%%%%%%%%%%%%%%%%%%%%%%%

A chance constraint requires a constraint to hold with a specified
probability.

A standard form is

\[
\boxed{
\mathbb{P}
\left(
g(x,\xi)\leq0
\right)
\geq
1-\varepsilon,
}
\]

where

\[
0<\varepsilon<1.
\]

Thus constraint violation is allowed only with probability at most
\(\varepsilon\).

%%%%%%%%%%%%%%%%%%%%%%%%%%%%%%%%%%%%%%%%%%%%%%%%%%%%%%%%%%%%
\subsection{Reliability Interpretation}
\label{subsec:chance-constraint-reliability}
%%%%%%%%%%%%%%%%%%%%%%%%%%%%%%%%%%%%%%%%%%%%%%%%%%%%%%%%%%%%

If

\[
\varepsilon=0.05,
\]

then

\[
\boxed{
\mathbb{P}
(
g(x,\xi)\leq0
)
\geq0.95.
}
\]

The constraint must therefore hold with at least \(95\%\) probability.

Chance constraints are useful when absolute robustness would be too
conservative.

%%%%%%%%%%%%%%%%%%%%%%%%%%%%%%%%%%%%%%%%%%%%%%%%%%%%%%%%%%%%
\subsection{Worked Example: Capacity Reliability}
\label{subsec:worked-chance-constraint}
%%%%%%%%%%%%%%%%%%%%%%%%%%%%%%%%%%%%%%%%%%%%%%%%%%%%%%%%%%%%

\begin{workedexamplebox}{Demand Coverage}

Suppose capacity \(x\) must cover random demand \(D\) with probability at
least \(0.99\).

The chance constraint is

\[
\boxed{
\mathbb{P}
(
D\leq x
)
\geq
0.99.
}
\]

If \(F_D\) is the demand cumulative distribution function, then

\[
F_D(x)\geq0.99.
\]

Thus a smallest feasible capacity is the \(0.99\)-quantile:

\begin{answerbox}
\[
\boxed{
x
\geq
F_D^{-1}(0.99)
}
\]

under the appropriate quantile definition.
\end{answerbox}

\end{workedexamplebox}

%%%%%%%%%%%%%%%%%%%%%%%%%%%%%%%%%%%%%%%%%%%%%%%%%%%%%%%%%%%%
\subsection{Gaussian Chance Constraint Example}
\label{subsec:gaussian-chance-constraint}
%%%%%%%%%%%%%%%%%%%%%%%%%%%%%%%%%%%%%%%%%%%%%%%%%%%%%%%%%%%%

Suppose

\[
D
\sim
N(\mu,\sigma^2).
\]

To require

\[
\mathbb{P}(D\leq x)
\geq
1-\varepsilon,
\]

one may write

\[
\boxed{
x
\geq
\mu
+
z_{1-\varepsilon}\sigma,
}
\]

where

\[
z_{1-\varepsilon}
\]

is the corresponding standard-normal quantile.

Thus some chance constraints admit deterministic reformulations.

%%%%%%%%%%%%%%%%%%%%%%%%%%%%%%%%%%%%%%%%%%%%%%%%%%%%%%%%%%%%
\subsection{Joint Chance Constraints}
\label{subsec:joint-chance-constraints}
%%%%%%%%%%%%%%%%%%%%%%%%%%%%%%%%%%%%%%%%%%%%%%%%%%%%%%%%%%%%

A joint chance constraint requires several inequalities to hold
simultaneously:

\[
\boxed{
\mathbb{P}
\left(
g_i(x,\xi)\leq0
\text{ for all }i
\right)
\geq
1-\varepsilon.
}
\]

Joint chance constraints are generally more difficult than individual
chance constraints.

Dependence among uncertain quantities becomes important.

%%%%%%%%%%%%%%%%%%%%%%%%%%%%%%%%%%%%%%%%%%%%%%%%%%%%%%%%%%%%
\subsection{Scenario Approximation of Chance Constraints}
\label{subsec:scenario-approximation-chance}
%%%%%%%%%%%%%%%%%%%%%%%%%%%%%%%%%%%%%%%%%%%%%%%%%%%%%%%%%%%%

One practical approximation draws samples

\[
\xi^{(1)},\ldots,\xi^{(N)}
\]

and enforces

\[
\boxed{
g
(
x,\xi^{(i)}
)
\leq0,
\qquad
i=1,\ldots,N.
}
\]

This replaces a probabilistic constraint with many deterministic sampled
constraints.

The statistical reliability of the resulting solution depends on sample
size and assumptions.

%%%%%%%%%%%%%%%%%%%%%%%%%%%%%%%%%%%%%%%%%%%%%%%%%%%%%%%%%%%%
\subsection{Risk Measures}
\label{subsec:risk-measures-stochastic}
%%%%%%%%%%%%%%%%%%%%%%%%%%%%%%%%%%%%%%%%%%%%%%%%%%%%%%%%%%%%

A risk measure maps a random loss \(L\) to a number:

\[
\boxed{
\rho(L).
}
\]

Possible risk measures include:

\[
\boxed{
\mathbb{E}[L],
}
\]

\[
\boxed{
\operatorname{Var}(L),
}
\]

\[
\boxed{
\operatorname{VaR}_{\alpha}(L),
}
\]

and

\[
\boxed{
\operatorname{CVaR}_{\alpha}(L).
}
\]

%%%%%%%%%%%%%%%%%%%%%%%%%%%%%%%%%%%%%%%%%%%%%%%%%%%%%%%%%%%%
\subsection{Mean--Variance Optimization}
\label{subsec:mean-variance-optimization}
%%%%%%%%%%%%%%%%%%%%%%%%%%%%%%%%%%%%%%%%%%%%%%%%%%%%%%%%%%%%

A mean--variance objective may balance expected cost and variance:

\[
\boxed{
\min_x
\left[
\mathbb{E}
[
C(x,\xi)
]
+
\lambda
\operatorname{Var}
[
C(x,\xi)
]
\right].
}
\]

The parameter

\[
\lambda\geq0
\]

controls risk aversion.

Larger \(\lambda\) penalizes variability more strongly.

%%%%%%%%%%%%%%%%%%%%%%%%%%%%%%%%%%%%%%%%%%%%%%%%%%%%%%%%%%%%
\subsection{Mean--Variance Portfolio Model}
\label{subsec:mean-variance-portfolio}
%%%%%%%%%%%%%%%%%%%%%%%%%%%%%%%%%%%%%%%%%%%%%%%%%%%%%%%%%%%%

Let

\[
w
\]

be portfolio weights,

\[
\mu
\]

expected returns, and

\[
\Sigma
\]

the covariance matrix.

A common model minimizes variance subject to a return target:

\[
\boxed{
\begin{aligned}
\min_w
\quad&
w^T\Sigma w
\\
\text{subject to}
\quad&
\mu^Tw
\geq
r_0,
\\
&
\mathbf{1}^Tw=1.
\end{aligned}
}
\]

When

\[
\Sigma\succeq0,
\]

this is a convex quadratic program.

%%%%%%%%%%%%%%%%%%%%%%%%%%%%%%%%%%%%%%%%%%%%%%%%%%%%%%%%%%%%
\subsection{Value-at-Risk}
\label{subsec:value-at-risk}
%%%%%%%%%%%%%%%%%%%%%%%%%%%%%%%%%%%%%%%%%%%%%%%%%%%%%%%%%%%%

For random loss \(L\), Value-at-Risk at level \(\alpha\) is a quantile of
the loss distribution.

A common definition is

\[
\boxed{
\operatorname{VaR}_{\alpha}(L)
=
\inf
\{
\ell:
\mathbb{P}(L\leq\ell)
\geq\alpha
\}.
}
\]

Thus VaR identifies a loss threshold not exceeded with probability at
least \(\alpha\).

%%%%%%%%%%%%%%%%%%%%%%%%%%%%%%%%%%%%%%%%%%%%%%%%%%%%%%%%%%%%
\subsection{Limitations of Value-at-Risk}
\label{subsec:limitations-var}
%%%%%%%%%%%%%%%%%%%%%%%%%%%%%%%%%%%%%%%%%%%%%%%%%%%%%%%%%%%%

VaR describes a quantile but does not directly describe the severity of
losses beyond that quantile.

Two loss distributions may have the same VaR while having very different
tail losses.

This motivates Conditional Value-at-Risk.

%%%%%%%%%%%%%%%%%%%%%%%%%%%%%%%%%%%%%%%%%%%%%%%%%%%%%%%%%%%%
\subsection{Conditional Value-at-Risk}
\label{subsec:conditional-value-at-risk}
%%%%%%%%%%%%%%%%%%%%%%%%%%%%%%%%%%%%%%%%%%%%%%%%%%%%%%%%%%%%

Conditional Value-at-Risk, abbreviated CVaR, measures average loss in the
upper tail beyond a specified confidence level.

For sufficiently regular continuous distributions, it may be interpreted
as

\[
\boxed{
\operatorname{CVaR}_{\alpha}(L)
=
\mathbb{E}
[
L
\mid
L\geq
\operatorname{VaR}_{\alpha}(L)
].
}
\]

For general distributions, a more careful optimization-based definition
is preferable.

%%%%%%%%%%%%%%%%%%%%%%%%%%%%%%%%%%%%%%%%%%%%%%%%%%%%%%%%%%%%
\subsection{Optimization Representation of CVaR}
\label{subsec:cvar-optimization-representation}
%%%%%%%%%%%%%%%%%%%%%%%%%%%%%%%%%%%%%%%%%%%%%%%%%%%%%%%%%%%%

For loss \(L(x,\xi)\), CVaR can be represented as

\[
\boxed{
\operatorname{CVaR}_{\alpha}
(
L(x,\xi)
)
=
\min_{\eta\in\mathbb{R}}
\left[
\eta
+
\frac{
1
}{
1-\alpha
}
\mathbb{E}
\left[
(
L(x,\xi)-\eta
)_+
\right]
\right].
}
\]

Here

\[
(u)_+
=
\max\{u,0\}.
\]

%%%%%%%%%%%%%%%%%%%%%%%%%%%%%%%%%%%%%%%%%%%%%%%%%%%%%%%%%%%%
\subsection{Scenario CVaR Formulation}
\label{subsec:scenario-cvar-formulation}
%%%%%%%%%%%%%%%%%%%%%%%%%%%%%%%%%%%%%%%%%%%%%%%%%%%%%%%%%%%%

Suppose losses under scenarios are

\[
L_\omega(x),
\]

with probabilities

\[
p_\omega.
\]

Introduce variables

\[
z_\omega\geq0
\]

such that

\[
\boxed{
z_\omega
\geq
L_\omega(x)-\eta.
}
\]

Then minimize

\[
\boxed{
\eta
+
\frac{
1
}{
1-\alpha
}
\sum_{\omega}
p_\omega z_\omega.
}
\]

This often produces a tractable convex formulation.

%%%%%%%%%%%%%%%%%%%%%%%%%%%%%%%%%%%%%%%%%%%%%%%%%%%%%%%%%%%%
\subsection{Worked Example: CVaR Auxiliary Variables}
\label{subsec:worked-cvar}
%%%%%%%%%%%%%%%%%%%%%%%%%%%%%%%%%%%%%%%%%%%%%%%%%%%%%%%%%%%%

\begin{workedexamplebox}{Tail-Loss Linearization}

Suppose scenario losses are

\[
L_1(x),
\quad
L_2(x),
\quad
L_3(x),
\]

with equal probabilities.

Introduce

\[
z_i
\geq
L_i(x)-\eta,
\]

and

\[
z_i\geq0.
\]

Then a sample CVaR objective is

\[
\boxed{
\eta
+
\frac{
1
}{
3(1-\alpha)
}
\sum_{i=1}^{3}
z_i.
}
\]

The positive-part terms have been represented using linear inequalities.

\end{workedexamplebox}

%%%%%%%%%%%%%%%%%%%%%%%%%%%%%%%%%%%%%%%%%%%%%%%%%%%%%%%%%%%%
\subsection{Risk-Averse Stochastic Optimization}
\label{subsec:risk-averse-stochastic-optimization}
%%%%%%%%%%%%%%%%%%%%%%%%%%%%%%%%%%%%%%%%%%%%%%%%%%%%%%%%%%%%

Instead of minimizing expected cost alone,

\[
\mathbb{E}[L],
\]

a risk-averse model may solve

\[
\boxed{
\min_x
\rho
(
L(x,\xi)
),
}
\]

where \(\rho\) penalizes undesirable uncertainty.

One may also combine expectation and tail risk:

\[
\boxed{
\min_x
\left[
\mathbb{E}[L]
+
\lambda
\operatorname{CVaR}_{\alpha}(L)
\right].
}
\]

%%%%%%%%%%%%%%%%%%%%%%%%%%%%%%%%%%%%%%%%%%%%%%%%%%%%%%%%%%%%
\subsection{Coherent Risk Measures Preview}
\label{subsec:coherent-risk-measures}
%%%%%%%%%%%%%%%%%%%%%%%%%%%%%%%%%%%%%%%%%%%%%%%%%%%%%%%%%%%%

A coherent risk measure satisfies properties such as:

\[
\boxed{
\text{monotonicity},
}
\]

\[
\boxed{
\text{subadditivity},
}
\]

\[
\boxed{
\text{positive homogeneity},
}
\]

and

\[
\boxed{
\text{translation invariance}.
}
\]

CVaR is an important example under standard formulations.

VaR is not generally coherent because it may fail subadditivity.

%%%%%%%%%%%%%%%%%%%%%%%%%%%%%%%%%%%%%%%%%%%%%%%%%%%%%%%%%%%%
\subsection{Robust Optimization}
\label{subsec:robust-optimization-stochastic}
%%%%%%%%%%%%%%%%%%%%%%%%%%%%%%%%%%%%%%%%%%%%%%%%%%%%%%%%%%%%

Robust optimization does not necessarily assign probabilities to
uncertain parameters.

Instead, assume

\[
\boxed{
\xi\in\mathcal{U},
}
\]

where \(\mathcal{U}\) is an uncertainty set.

A robust constraint requires

\[
\boxed{
g(x,\xi)\leq0
\qquad
\text{for every }
\xi\in\mathcal{U}.
}
\]

%%%%%%%%%%%%%%%%%%%%%%%%%%%%%%%%%%%%%%%%%%%%%%%%%%%%%%%%%%%%
\subsection{Stochastic Versus Robust Optimization}
\label{subsec:stochastic-versus-robust}
%%%%%%%%%%%%%%%%%%%%%%%%%%%%%%%%%%%%%%%%%%%%%%%%%%%%%%%%%%%%

Stochastic optimization uses:

\[
\boxed{
\text{probability distributions}.
}
\]

Robust optimization uses:

\[
\boxed{
\text{uncertainty sets}.
}
\]

A stochastic model may optimize average or probabilistic performance.

A robust model protects against all uncertainty realizations in the chosen
set.

%%%%%%%%%%%%%%%%%%%%%%%%%%%%%%%%%%%%%%%%%%%%%%%%%%%%%%%%%%%%
\subsection{Conservatism in Robust Optimization}
\label{subsec:robust-conservatism}
%%%%%%%%%%%%%%%%%%%%%%%%%%%%%%%%%%%%%%%%%%%%%%%%%%%%%%%%%%%%

If \(\mathcal{U}\) is very large, robust solutions may be overly
conservative.

If \(\mathcal{U}\) is too small, important uncertainty may be ignored.

Thus uncertainty-set design controls the tradeoff between:

\[
\boxed{
\text{protection}
}
\]

and

\[
\boxed{
\text{performance}.
}
\]

%%%%%%%%%%%%%%%%%%%%%%%%%%%%%%%%%%%%%%%%%%%%%%%%%%%%%%%%%%%%
\subsection{Distributionally Robust Optimization}
\label{subsec:distributionally-robust-optimization}
%%%%%%%%%%%%%%%%%%%%%%%%%%%%%%%%%%%%%%%%%%%%%%%%%%%%%%%%%%%%

Distributionally robust optimization, abbreviated DRO, assumes the exact
probability distribution is unknown.

Instead,

\[
\boxed{
P
\in
\mathcal{P},
}
\]

where \(\mathcal{P}\) is an ambiguity set of distributions.

A typical problem is

\[
\boxed{
\min_x
\sup_{P\in\mathcal{P}}
\mathbb{E}_{P}
[
F(x,\xi)
].
}
\]

%%%%%%%%%%%%%%%%%%%%%%%%%%%%%%%%%%%%%%%%%%%%%%%%%%%%%%%%%%%%
\subsection{Ambiguity Sets}
\label{subsec:ambiguity-sets}
%%%%%%%%%%%%%%%%%%%%%%%%%%%%%%%%%%%%%%%%%%%%%%%%%%%%%%%%%%%%

An ambiguity set may contain distributions satisfying information such as:

\[
\boxed{
\text{known moments},
}
\]

\[
\boxed{
\text{support restrictions},
}
\]

\[
\boxed{
\text{distance from an empirical distribution},
}
\]

or

\[
\boxed{
\text{confidence-region constraints}.
}
\]

DRO protects against distributional misspecification.

%%%%%%%%%%%%%%%%%%%%%%%%%%%%%%%%%%%%%%%%%%%%%%%%%%%%%%%%%%%%
\subsection{Aleatoric and Epistemic Uncertainty}
\label{subsec:aleatoric-epistemic-uncertainty}
%%%%%%%%%%%%%%%%%%%%%%%%%%%%%%%%%%%%%%%%%%%%%%%%%%%%%%%%%%%%

Aleatoric uncertainty refers to intrinsic randomness in the system.

Examples include:

\[
\boxed{
\text{random demand}
}
\]

or

\[
\boxed{
\text{random failures}.
}
\]

Epistemic uncertainty refers to imperfect knowledge about the model or its
parameters.

Distributionally robust methods are particularly relevant when the
probability law itself is uncertain.

%%%%%%%%%%%%%%%%%%%%%%%%%%%%%%%%%%%%%%%%%%%%%%%%%%%%%%%%%%%%
\subsection{Decomposition of Two-Stage Problems}
\label{subsec:decomposition-two-stage-stochastic}
%%%%%%%%%%%%%%%%%%%%%%%%%%%%%%%%%%%%%%%%%%%%%%%%%%%%%%%%%%%%

A large finite-scenario two-stage model can have many scenario-specific
variables.

However, scenarios are often coupled only through the first-stage
decision.

This special structure allows decomposition.

%%%%%%%%%%%%%%%%%%%%%%%%%%%%%%%%%%%%%%%%%%%%%%%%%%%%%%%%%%%%
\subsection{Benders Decomposition Connection}
\label{subsec:benders-stochastic}
%%%%%%%%%%%%%%%%%%%%%%%%%%%%%%%%%%%%%%%%%%%%%%%%%%%%%%%%%%%%

Fix a first-stage decision \(x\).

Each scenario recourse problem can then be solved independently.

Dual information from the recourse problems generates cuts for a master
problem in \(x\).

In stochastic programming, this leads to the L-shaped method.

%%%%%%%%%%%%%%%%%%%%%%%%%%%%%%%%%%%%%%%%%%%%%%%%%%%%%%%%%%%%
\subsection{L-Shaped Method}
\label{subsec:l-shaped-method}
%%%%%%%%%%%%%%%%%%%%%%%%%%%%%%%%%%%%%%%%%%%%%%%%%%%%%%%%%%%%

For suitable two-stage stochastic linear programs, the L-shaped method
iteratively:

\begin{enumerate}
    \item solves a first-stage master problem;
    \item evaluates scenario recourse subproblems;
    \item generates feasibility cuts if recourse is infeasible;
    \item generates optimality cuts estimating expected recourse cost;
    \item resolves the master problem.
\end{enumerate}

It is a specialized Benders-decomposition method.

%%%%%%%%%%%%%%%%%%%%%%%%%%%%%%%%%%%%%%%%%%%%%%%%%%%%%%%%%%%%
\subsection{Recourse Approximation}
\label{subsec:recourse-approximation}
%%%%%%%%%%%%%%%%%%%%%%%%%%%%%%%%%%%%%%%%%%%%%%%%%%%%%%%%%%%%

The expected recourse function

\[
\boxed{
Q(x)
=
\mathbb{E}
[
Q(x,\xi)
]
}
\]

may be difficult to evaluate directly.

Decomposition methods construct lower approximations using supporting
cuts.

The approximation becomes tighter as more recourse information is added.

%%%%%%%%%%%%%%%%%%%%%%%%%%%%%%%%%%%%%%%%%%%%%%%%%%%%%%%%%%%%
\subsection{Stochastic Dual Dynamic Programming Preview}
\label{subsec:stochastic-dual-dynamic-programming}
%%%%%%%%%%%%%%%%%%%%%%%%%%%%%%%%%%%%%%%%%%%%%%%%%%%%%%%%%%%%

Stochastic Dual Dynamic Programming, abbreviated SDDP, is used for large
multistage stochastic linear or convex problems with suitable structure.

It approximates future value functions using collections of cuts.

A typical iteration includes:

\[
\boxed{
\text{forward simulation},
}
\]

and

\[
\boxed{
\text{backward cut generation}.
}
\]

%%%%%%%%%%%%%%%%%%%%%%%%%%%%%%%%%%%%%%%%%%%%%%%%%%%%%%%%%%%%
\subsection{Forward Pass in SDDP}
\label{subsec:sddp-forward-pass}
%%%%%%%%%%%%%%%%%%%%%%%%%%%%%%%%%%%%%%%%%%%%%%%%%%%%%%%%%%%%

The forward pass samples uncertainty and generates a candidate policy
trajectory through the stages.

It identifies states where improved value-function approximations may be
useful.

%%%%%%%%%%%%%%%%%%%%%%%%%%%%%%%%%%%%%%%%%%%%%%%%%%%%%%%%%%%%
\subsection{Backward Pass in SDDP}
\label{subsec:sddp-backward-pass}
%%%%%%%%%%%%%%%%%%%%%%%%%%%%%%%%%%%%%%%%%%%%%%%%%%%%%%%%%%%%

The backward pass solves stage subproblems and uses dual information to
construct new supporting cuts.

These cuts improve the approximations of future cost functions.

Repeated iterations gradually strengthen the policy.

%%%%%%%%%%%%%%%%%%%%%%%%%%%%%%%%%%%%%%%%%%%%%%%%%%%%%%%%%%%%
\subsection{Stochastic Programming and Convexity}
\label{subsec:stochastic-programming-convexity}
%%%%%%%%%%%%%%%%%%%%%%%%%%%%%%%%%%%%%%%%%%%%%%%%%%%%%%%%%%%%

If

\[
F(x,\xi)
\]

is convex in \(x\) for every \(\xi\), then under appropriate integrability
conditions,

\[
\boxed{
f(x)
=
\mathbb{E}
[
F(x,\xi)
]
}
\]

is convex.

Thus expectations preserve convexity.

This is one reason many stochastic programs remain tractable.

%%%%%%%%%%%%%%%%%%%%%%%%%%%%%%%%%%%%%%%%%%%%%%%%%%%%%%%%%%%%
\subsection{Proof that Expectation Preserves Convexity}
\label{subsec:expectation-preserves-convexity}
%%%%%%%%%%%%%%%%%%%%%%%%%%%%%%%%%%%%%%%%%%%%%%%%%%%%%%%%%%%%

Suppose \(F(\cdot,\xi)\) is convex for every \(\xi\).

Then

\[
F
(
\theta x+(1-\theta)y,
\xi
)
\leq
\theta F(x,\xi)
+
(1-\theta)F(y,\xi).
\]

Take expectations:

\[
\mathbb{E}
[
F
(
\theta x+(1-\theta)y,
\xi
)
]
\leq
\theta
\mathbb{E}
[
F(x,\xi)
]
+
(1-\theta)
\mathbb{E}
[
F(y,\xi)
].
\]

Therefore,

\[
\boxed{
\mathbb{E}[F(x,\xi)]
\text{ is convex in }x.
}
\]

%%%%%%%%%%%%%%%%%%%%%%%%%%%%%%%%%%%%%%%%%%%%%%%%%%%%%%%%%%%%
\subsection{Interchanging Gradient and Expectation}
\label{subsec:gradient-expectation-interchange}
%%%%%%%%%%%%%%%%%%%%%%%%%%%%%%%%%%%%%%%%%%%%%%%%%%%%%%%%%%%%

Under suitable differentiability, integrability, and domination
conditions,

\[
\boxed{
\nabla
\mathbb{E}
[
F(x,\xi)
]
=
\mathbb{E}
[
\nabla_xF(x,\xi)
].
}
\]

This identity justifies stochastic gradient methods.

It should not be assumed without appropriate regularity.

%%%%%%%%%%%%%%%%%%%%%%%%%%%%%%%%%%%%%%%%%%%%%%%%%%%%%%%%%%%%
\subsection{Jensen's Inequality Connection}
\label{subsec:jensen-stochastic-optimization}
%%%%%%%%%%%%%%%%%%%%%%%%%%%%%%%%%%%%%%%%%%%%%%%%%%%%%%%%%%%%

For convex \(f\),

\[
\boxed{
f
(
\mathbb{E}[X]
)
\leq
\mathbb{E}
[
f(X)
].
}
\]

Therefore replacing uncertainty by its mean can systematically
underestimate the expected value of a convex cost.

This explains one way in which the expected-value problem can differ from
the true stochastic problem.

%%%%%%%%%%%%%%%%%%%%%%%%%%%%%%%%%%%%%%%%%%%%%%%%%%%%%%%%%%%%
\subsection{Worked Example: Jensen's Inequality}
\label{subsec:worked-jensen-stochastic}
%%%%%%%%%%%%%%%%%%%%%%%%%%%%%%%%%%%%%%%%%%%%%%%%%%%%%%%%%%%%

\begin{workedexamplebox}{Cost of Random Demand}

Let

\[
f(D)=D^2.
\]

Suppose

\[
\mathbb{E}[D]=10
\]

and

\[
\operatorname{Var}(D)=4.
\]

Then

\[
\mathbb{E}[D^2]
=
\operatorname{Var}(D)
+
(\mathbb{E}[D])^2.
\]

Thus,

\[
\mathbb{E}[D^2]
=
4+100
=
104.
\]

But

\[
f(\mathbb{E}[D])
=
10^2
=
100.
\]

Therefore,

\begin{answerbox}
\[
\boxed{
f(\mathbb{E}[D])
=
100
<
104
=
\mathbb{E}[f(D)].
}
\]
\end{answerbox}

\end{workedexamplebox}

%%%%%%%%%%%%%%%%%%%%%%%%%%%%%%%%%%%%%%%%%%%%%%%%%%%%%%%%%%%%
\subsection{Decision Rules}
\label{subsec:stochastic-decision-rules}
%%%%%%%%%%%%%%%%%%%%%%%%%%%%%%%%%%%%%%%%%%%%%%%%%%%%%%%%%%%%

In multistage problems, a future decision is a function of observed
information.

Thus instead of optimizing a finite vector alone, one may optimize a
policy:

\[
\boxed{
x_t
=
\pi_t
(
\xi_1,\ldots,\xi_t
).
}
\]

The full space of possible functions can be enormous.

%%%%%%%%%%%%%%%%%%%%%%%%%%%%%%%%%%%%%%%%%%%%%%%%%%%%%%%%%%%%
\subsection{Affine Decision Rules Preview}
\label{subsec:affine-decision-rules}
%%%%%%%%%%%%%%%%%%%%%%%%%%%%%%%%%%%%%%%%%%%%%%%%%%%%%%%%%%%%

A tractable approximation restricts decisions to affine functions of
observed uncertainty:

\[
\boxed{
x_t
=
a_t
+
B_t
\boldsymbol{\xi}_{1:t}.
}
\]

This reduces an infinite-dimensional policy problem to a finite number of
parameters.

The resulting policy is approximate but computationally manageable.

%%%%%%%%%%%%%%%%%%%%%%%%%%%%%%%%%%%%%%%%%%%%%%%%%%%%%%%%%%%%
\subsection{Rolling-Horizon Stochastic Optimization}
\label{subsec:rolling-horizon-stochastic}
%%%%%%%%%%%%%%%%%%%%%%%%%%%%%%%%%%%%%%%%%%%%%%%%%%%%%%%%%%%%

A model can be repeatedly re-solved as new information becomes available.

At each stage:

\begin{enumerate}
    \item update forecasts or distributions;
    \item solve the optimization model;
    \item implement the first decision;
    \item observe new information;
    \item repeat.
\end{enumerate}

This creates an adaptive rolling-horizon policy.

%%%%%%%%%%%%%%%%%%%%%%%%%%%%%%%%%%%%%%%%%%%%%%%%%%%%%%%%%%%%
\subsection{Forecast Error and Optimization}
\label{subsec:forecast-error-optimization}
%%%%%%%%%%%%%%%%%%%%%%%%%%%%%%%%%%%%%%%%%%%%%%%%%%%%%%%%%%%%

A forecast provides an estimate such as

\[
\widehat{D}.
\]

A stochastic model instead represents forecast uncertainty through a
distribution or scenario set.

Thus:

\[
\boxed{
\text{forecast}
}
\]

and

\[
\boxed{
\text{uncertainty model}
}
\]

are distinct components.

Optimization quality depends on both.

%%%%%%%%%%%%%%%%%%%%%%%%%%%%%%%%%%%%%%%%%%%%%%%%%%%%%%%%%%%%
\subsection{Probability Estimation Error}
\label{subsec:probability-estimation-error}
%%%%%%%%%%%%%%%%%%%%%%%%%%%%%%%%%%%%%%%%%%%%%%%%%%%%%%%%%%%%

Scenario probabilities are rarely known exactly.

They may be estimated from finite data.

Therefore the optimization problem itself can be sensitive to statistical
estimation error.

This motivates:

\[
\boxed{
\text{out-of-sample testing},
}
\]

\[
\boxed{
\text{regularization},
}
\]

and

\[
\boxed{
\text{distributionally robust optimization}.
}
\]

%%%%%%%%%%%%%%%%%%%%%%%%%%%%%%%%%%%%%%%%%%%%%%%%%%%%%%%%%%%%
\subsection{Rare Events}
\label{subsec:rare-events-stochastic-optimization}
%%%%%%%%%%%%%%%%%%%%%%%%%%%%%%%%%%%%%%%%%%%%%%%%%%%%%%%%%%%%

Rare but severe outcomes may contribute little to ordinary sample averages
unless very large samples are used.

Yet they may dominate risk.

Examples include:

\[
\boxed{
\text{extreme financial losses},
}
\]

\[
\boxed{
\text{system failures},
}
\]

and

\[
\boxed{
\text{natural disasters}.
}
\]

Tail-risk measures and importance-sampling ideas may therefore be useful.

%%%%%%%%%%%%%%%%%%%%%%%%%%%%%%%%%%%%%%%%%%%%%%%%%%%%%%%%%%%%
\subsection{Importance Sampling Preview}
\label{subsec:importance-sampling-stochastic-optimization}
%%%%%%%%%%%%%%%%%%%%%%%%%%%%%%%%%%%%%%%%%%%%%%%%%%%%%%%%%%%%

Importance sampling changes the sampling distribution to generate
important rare events more frequently.

Statistical weights correct for the altered distribution.

This can reduce variance when estimating rare-event probabilities or tail
costs.

The method is developed more deeply with Monte Carlo techniques in
Chapter 11.

%%%%%%%%%%%%%%%%%%%%%%%%%%%%%%%%%%%%%%%%%%%%%%%%%%%%%%%%%%%%
\subsection{Stochastic Optimization Workflow}
\label{subsec:stochastic-optimization-workflow}
%%%%%%%%%%%%%%%%%%%%%%%%%%%%%%%%%%%%%%%%%%%%%%%%%%%%%%%%%%%%

A practical workflow is:

\begin{enumerate}
    \item identify uncertain quantities;
    \item determine when each uncertain quantity becomes observable;
    \item define here-and-now and recourse decisions;
    \item choose a probability model, scenario model, or ambiguity set;
    \item formulate expected cost, risk, or chance constraints;
    \item enforce nonanticipativity;
    \item identify whether the problem is static, two-stage, or multistage;
    \item exploit convexity where available;
    \item determine whether exact expectations are tractable;
    \item use scenarios or sample-average approximation when needed;
    \item choose decomposition for large scenario models;
    \item compute or estimate statistical error;
    \item evaluate solutions out of sample;
    \item examine tail risk and rare events;
    \item compare stochastic and deterministic expected-value solutions;
    \item calculate VSS or EVPI when useful;
    \item perform sensitivity analysis on distributions and probabilities;
    \item interpret the final policy rather than only the objective value.
\end{enumerate}

%%%%%%%%%%%%%%%%%%%%%%%%%%%%%%%%%%%%%%%%%%%%%%%%%%%%%%%%%%%%
\subsection{Choosing a Stochastic Optimization Method}
\label{subsec:choosing-stochastic-method}
%%%%%%%%%%%%%%%%%%%%%%%%%%%%%%%%%%%%%%%%%%%%%%%%%%%%%%%%%%%%

For finite scenario two-stage linear models, consider:

\[
\boxed{
\text{deterministic equivalent}
}
\]

or

\[
\boxed{
\text{L-shaped decomposition}.
}
\]

For large expectations estimated from data, consider:

\[
\boxed{
\text{sample-average approximation}.
}
\]

For very large smooth expected objectives, consider:

\[
\boxed{
\text{stochastic gradient methods}.
}
\]

For multistage problems with manageable states, consider:

\[
\boxed{
\text{stochastic dynamic programming}.
}
\]

For large multistage convex problems, consider methods such as:

\[
\boxed{
\text{SDDP}.
}
\]

For probabilistic reliability requirements, consider:

\[
\boxed{
\text{chance constraints}.
}
\]

For uncertain probability distributions, consider:

\[
\boxed{
\text{distributionally robust optimization}.
}
\]

%%%%%%%%%%%%%%%%%%%%%%%%%%%%%%%%%%%%%%%%%%%%%%%%%%%%%%%%%%%%
\subsection{Common Errors}
\label{subsec:stochastic-optimization-common-errors}
%%%%%%%%%%%%%%%%%%%%%%%%%%%%%%%%%%%%%%%%%%%%%%%%%%%%%%%%%%%%

\subsubsection{Replacing Random Data by Its Mean Automatically}

In general,

\[
\boxed{
\min_x
F
(
x,\mathbb{E}[\xi]
)
}
\]

is not equivalent to

\[
\boxed{
\min_x
\mathbb{E}
[
F(x,\xi)
].
}
\]

Nonlinearity and recourse can make the solutions very different.

%%%%%%%%%%%%%%%%%%%%%%%%%%%%%%%%%%%%%%%%%%%%%%%%%%%%%%%%%%%%
\subsubsection{Ignoring Information Timing}

A decision may depend only on uncertainty already observed.

Using future information violates nonanticipativity.

%%%%%%%%%%%%%%%%%%%%%%%%%%%%%%%%%%%%%%%%%%%%%%%%%%%%%%%%%%%%
\subsubsection{Letting First-Stage Decisions Vary by Scenario}

If a decision is made before uncertainty occurs, it must be identical
across scenarios.

Creating separate first-stage values without nonanticipativity gives the
model impossible foresight.

%%%%%%%%%%%%%%%%%%%%%%%%%%%%%%%%%%%%%%%%%%%%%%%%%%%%%%%%%%%%
\subsubsection{Assuming Expected Cost Controls Tail Risk}

A small expected cost does not guarantee small worst-case or tail losses.

Risk measures should match the application.

%%%%%%%%%%%%%%%%%%%%%%%%%%%%%%%%%%%%%%%%%%%%%%%%%%%%%%%%%%%%
\subsubsection{Confusing VaR and CVaR}

VaR is fundamentally a quantile.

CVaR describes upper-tail loss through an averaging or optimization
construction.

They are not interchangeable.

%%%%%%%%%%%%%%%%%%%%%%%%%%%%%%%%%%%%%%%%%%%%%%%%%%%%%%%%%%%%
\subsubsection{Treating a Sample Average as the Exact Expectation}

An SAA objective contains sampling error.

Different random samples may produce different solutions.

%%%%%%%%%%%%%%%%%%%%%%%%%%%%%%%%%%%%%%%%%%%%%%%%%%%%%%%%%%%%
\subsubsection{Evaluating on the Same Scenarios Used for Optimization}

This can make a solution appear better than it performs on new data.

Independent out-of-sample evaluation is important.

%%%%%%%%%%%%%%%%%%%%%%%%%%%%%%%%%%%%%%%%%%%%%%%%%%%%%%%%%%%%
\subsubsection{Using Too Few Scenarios}

A small scenario set may miss important regions of the uncertainty
distribution.

The resulting optimization model can produce unstable decisions.

%%%%%%%%%%%%%%%%%%%%%%%%%%%%%%%%%%%%%%%%%%%%%%%%%%%%%%%%%%%%
\subsubsection{Ignoring Scenario Probabilities}

For unequal probabilities,

\[
\boxed{
\mathbb{E}[F]
=
\sum_\omega
p_\omega F_\omega.
}
\]

Using an unweighted average changes the model.

%%%%%%%%%%%%%%%%%%%%%%%%%%%%%%%%%%%%%%%%%%%%%%%%%%%%%%%%%%%%
\subsubsection{Confusing Stochastic and Robust Optimization}

Stochastic optimization uses probability information.

Robust optimization protects against all uncertainty in a chosen set.

The resulting decisions can be very different.

%%%%%%%%%%%%%%%%%%%%%%%%%%%%%%%%%%%%%%%%%%%%%%%%%%%%%%%%%%%%
\subsubsection{Assuming Chance Constraints Guarantee Zero Violations}

A chance constraint permits violation probability up to \(\varepsilon\).

It is not a hard constraint for every realization.

%%%%%%%%%%%%%%%%%%%%%%%%%%%%%%%%%%%%%%%%%%%%%%%%%%%%%%%%%%%%
\subsubsection{Ignoring Dependence in Joint Chance Constraints}

Individual reliability levels do not automatically determine joint
reliability.

Dependence among uncertain constraints matters.

%%%%%%%%%%%%%%%%%%%%%%%%%%%%%%%%%%%%%%%%%%%%%%%%%%%%%%%%%%%%
\subsubsection{Assuming Stochastic Gradients Equal Exact Gradients}

Individual stochastic gradients are noisy estimates.

Only their expectation may equal the true gradient under suitable
assumptions.

%%%%%%%%%%%%%%%%%%%%%%%%%%%%%%%%%%%%%%%%%%%%%%%%%%%%%%%%%%%%
\subsubsection{Using a Large Constant SGD Step Size Indefinitely}

Persistent gradient noise can cause continued fluctuations.

Step-size design affects convergence.

%%%%%%%%%%%%%%%%%%%%%%%%%%%%%%%%%%%%%%%%%%%%%%%%%%%%%%%%%%%%
\subsubsection{Interchanging Expectation and Differentiation Without Conditions}

The identity

\[
\nabla
\mathbb{E}[F]
=
\mathbb{E}[\nabla F]
\]

requires regularity assumptions.

It is not purely algebraic.

%%%%%%%%%%%%%%%%%%%%%%%%%%%%%%%%%%%%%%%%%%%%%%%%%%%%%%%%%%%%
\subsubsection{Interpreting EVPI with the Wrong Sign}

For standard minimization notation,

\[
\boxed{
EVPI=RP-WS\geq0.
}
\]

The sign changes if a different maximization convention is used.

%%%%%%%%%%%%%%%%%%%%%%%%%%%%%%%%%%%%%%%%%%%%%%%%%%%%%%%%%%%%
\subsubsection{Interpreting VSS with the Wrong Sign}

For standard minimization notation,

\[
\boxed{
VSS=EEV-RP\geq0.
}
\]

The definition must be consistent with the optimization direction.

%%%%%%%%%%%%%%%%%%%%%%%%%%%%%%%%%%%%%%%%%%%%%%%%%%%%%%%%%%%%
\subsection{Applications}
\label{subsec:stochastic-optimization-applications}
%%%%%%%%%%%%%%%%%%%%%%%%%%%%%%%%%%%%%%%%%%%%%%%%%%%%%%%%%%%%

\begin{applicationbox}

\textbf{Supply Chains.}
Production, inventory, transportation, and capacity decisions are made
before uncertain demand and disruptions are fully known.

\medskip

\textbf{Energy Systems.}
Generation planning must account for uncertain demand, renewable output,
fuel prices, and outages.

\medskip

\textbf{Finance.}
Portfolio and asset-allocation models balance uncertain returns with risk
measures such as variance and CVaR.

\medskip

\textbf{Transportation.}
Routing, fleet allocation, and scheduling may depend on uncertain travel
times and demand.

\medskip

\textbf{Inventory.}
Order quantities must be chosen while future demand remains uncertain.

\medskip

\textbf{Machine Learning.}
Training objectives are commonly expectations over data distributions and
are optimized using stochastic gradients.

\medskip

\textbf{Healthcare.}
Staffing, capacity, and resource allocation may need to account for
uncertain patient arrivals and treatment demand.

\medskip

\textbf{Control.}
Stochastic control optimizes feedback policies for systems affected by
random disturbances.

\end{applicationbox}

%%%%%%%%%%%%%%%%%%%%%%%%%%%%%%%%%%%%%%%%%%%%%%%%%%%%%%%%%%%%
\subsection{Exercises}
\label{subsec:stochastic-optimization-exercises}
%%%%%%%%%%%%%%%%%%%%%%%%%%%%%%%%%%%%%%%%%%%%%%%%%%%%%%%%%%%%

\begin{exercise}[1]
Define stochastic optimization.
\end{exercise}

\begin{exercise}[1]
Define a here-and-now decision.
\end{exercise}

\begin{exercise}[1]
Define a recourse decision.
\end{exercise}

\begin{exercise}[1]
Define a scenario.
\end{exercise}

\begin{exercise}[1]
Define nonanticipativity.
\end{exercise}

\begin{exercise}[1]
Define sample-average approximation.
\end{exercise}

\begin{exercise}[1]
Define a chance constraint.
\end{exercise}

\begin{exercise}[1]
Define Value-at-Risk conceptually.
\end{exercise}

\begin{exercise}[1]
Define Conditional Value-at-Risk conceptually.
\end{exercise}

\begin{exercise}[1]
Define distributionally robust optimization.
\end{exercise}

\begin{exercise}[2]
Suppose a cost is \(20\) with probability \(0.7\) and \(50\) with
probability \(0.3\). Compute its expected value.
\end{exercise}

\begin{exercise}[2]
Write the finite-scenario expectation for three scenarios with
probabilities \(0.2\), \(0.5\), and \(0.3\).
\end{exercise}

\begin{exercise}[2]
Explain the difference between first-stage and second-stage decisions.
\end{exercise}

\begin{exercise}[2]
Give an example of a recourse action in transportation.
\end{exercise}

\begin{exercise}[2]
Explain why first-stage variables must satisfy nonanticipativity.
\end{exercise}

\begin{exercise}[2]
Write a chance constraint requiring at least \(95\%\) service reliability.
\end{exercise}

\begin{exercise}[2]
If demand has distribution function \(F_D\), write a capacity requirement
that covers demand with probability at least \(0.9\).
\end{exercise}

\begin{exercise}[2]
Compute the sample average of costs

\[
8,\ 11,\ 13,\ 10,\ 8.
\]
\end{exercise}

\begin{exercise}[2]
Take one SGD step from \(x_0=5\) using stochastic gradient \(4\) and step
size \(0.2\).
\end{exercise}

\begin{exercise}[3]
Construct a two-stage capacity-expansion model with emergency recourse.
\end{exercise}

\begin{exercise}[3]
Write the deterministic equivalent of a two-scenario stochastic linear
program.
\end{exercise}

\begin{exercise}[3]
Explain how the deterministic equivalent changes as the number of
scenarios increases.
\end{exercise}

\begin{exercise}[3]
Construct a three-stage scenario tree with two branches at each uncertain
stage.
\end{exercise}

\begin{exercise}[3]
Identify which decisions must be equal under nonanticipativity in the
previous tree.
\end{exercise}

\begin{exercise}[3]
Suppose

\[
RP=85,
\qquad
WS=70.
\]

Compute EVPI.
\end{exercise}

\begin{exercise}[3]
Suppose

\[
EEV=96,
\qquad
RP=85.
\]

Compute VSS.
\end{exercise}

\begin{exercise}[3]
Explain the practical meaning of a large VSS.
\end{exercise}

\begin{exercise}[3]
Explain the practical meaning of a large EVPI.
\end{exercise}

\begin{exercise}[3]
Construct an SAA problem using \(N\) sampled scenarios.
\end{exercise}

\begin{exercise}[3]
Explain why SAA solutions should be tested out of sample.
\end{exercise}

\begin{exercise}[3]
Show that the sample mean is an unbiased estimator of an expected cost
under independent identically distributed sampling.
\end{exercise}

\begin{exercise}[3]
Explain the variance tradeoff between single-sample SGD and minibatch SGD.
\end{exercise}

\begin{exercise}[3]
Verify whether

\[
\alpha_k=\frac1k
\]

satisfies

\[
\sum_k\alpha_k=\infty
\]

and

\[
\sum_k\alpha_k^2<\infty.
\]
\end{exercise}

\begin{exercise}[3]
Write a mean--variance optimization objective.
\end{exercise}

\begin{exercise}[3]
Formulate a minimum-variance portfolio problem with a target expected
return.
\end{exercise}

\begin{exercise}[3]
For a loss distribution, identify the difference between VaR and CVaR.
\end{exercise}

\begin{exercise}[3]
Write the auxiliary-variable representation of CVaR.
\end{exercise}

\begin{exercise}[3]
Construct a finite-scenario CVaR optimization model.
\end{exercise}

\begin{exercise}[3]
Explain why CVaR focuses more directly on tail losses than expectation.
\end{exercise}

\begin{exercise}[3]
Compare a chance constraint with a robust constraint.
\end{exercise}

\begin{exercise}[3]
Suppose

\[
D\sim N(100,15^2).
\]

Write the capacity requirement for \(99\%\) service using the standard
normal quantile \(z_{0.99}\).
\end{exercise}

\begin{exercise}[3]
Explain why joint chance constraints require dependence information.
\end{exercise}

\begin{exercise}[3]
Show that expectation preserves convexity.
\end{exercise}

\begin{exercise}[3]
Use Jensen's inequality to explain why replacing a random input by its mean
may underestimate convex expected cost.
\end{exercise}

\begin{exercise}[3]
Formulate a distributionally robust expected-cost problem.
\end{exercise}

\begin{exercise}[3]
Give examples of possible ambiguity-set information.
\end{exercise}

\begin{exercise}[3]
Explain the difference between aleatoric and epistemic uncertainty.
\end{exercise}

\begin{exercise}[3]
Describe the major steps of the L-shaped method.
\end{exercise}

\begin{exercise}[3]
Explain why scenario recourse problems can often be solved independently
once the first-stage decision is fixed.
\end{exercise}

\begin{exercise}[3]
Explain the basic forward and backward passes of SDDP.
\end{exercise}

\begin{exercise}[4]
Study existence theory for two-stage stochastic programs.
\end{exercise}

\begin{exercise}[4]
Study convexity and continuity of recourse functions.
\end{exercise}

\begin{exercise}[4]
Derive the deterministic equivalent of a multistage scenario-tree model.
\end{exercise}

\begin{exercise}[4]
Study relatively complete and complete recourse.
\end{exercise}

\begin{exercise}[4]
Derive L-shaped optimality cuts from dual recourse solutions.
\end{exercise}

\begin{exercise}[4]
Study feasibility cuts for stochastic programs without relatively complete
recourse.
\end{exercise}

\begin{exercise}[4]
Study convergence theory for sample-average approximation.
\end{exercise}

\begin{exercise}[4]
Construct statistical confidence bounds for an SAA optimal value.
\end{exercise}

\begin{exercise}[4]
Study Robbins--Monro stochastic approximation.
\end{exercise}

\begin{exercise}[4]
Prove convergence results for stochastic gradient descent under convexity.
\end{exercise}

\begin{exercise}[4]
Study strongly convex stochastic-gradient convergence rates.
\end{exercise}

\begin{exercise}[4]
Study Polyak--Ruppert averaging.
\end{exercise}

\begin{exercise}[4]
Investigate stochastic variance-reduction methods.
\end{exercise}

\begin{exercise}[4]
Study deterministic reformulations of Gaussian chance constraints.
\end{exercise}

\begin{exercise}[4]
Study scenario-approach guarantees for chance-constrained optimization.
\end{exercise}

\begin{exercise}[4]
Study coherent risk measures.
\end{exercise}

\begin{exercise}[4]
Derive the optimization representation of CVaR.
\end{exercise}

\begin{exercise}[4]
Study time-consistent dynamic risk measures.
\end{exercise}

\begin{exercise}[4]
Study robust counterparts of linear programs.
\end{exercise}

\begin{exercise}[4]
Study ellipsoidal uncertainty sets and robust optimization.
\end{exercise}

\begin{exercise}[4]
Study Wasserstein ambiguity sets in distributionally robust optimization.
\end{exercise}

\begin{exercise}[4]
Study moment-based distributionally robust optimization.
\end{exercise}

\begin{exercise}[4]
Study affine decision rules for multistage uncertainty.
\end{exercise}

\begin{exercise}[4]
Study stochastic dual dynamic programming.
\end{exercise}

\begin{exercise}[4]
Investigate rare-event simulation for stochastic optimization.
\end{exercise}

%%%%%%%%%%%%%%%%%%%%%%%%%%%%%%%%%%%%%%%%%%%%%%%%%%%%%%%%%%%%
\subsection{Conceptual Summary}
\label{subsec:stochastic-optimization-summary}
%%%%%%%%%%%%%%%%%%%%%%%%%%%%%%%%%%%%%%%%%%%%%%%%%%%%%%%%%%%%

Stochastic optimization incorporates uncertain quantities directly into
optimization models.

A basic expected-value problem is

\[
\boxed{
\min_x
\mathbb{E}
[
F(x,\xi)
].
}
\]

A two-stage model separates:

\[
\boxed{
\text{here-and-now decisions}
}
\]

from

\[
\boxed{
\text{recourse decisions}.
}
\]

A finite-scenario objective takes form

\[
\boxed{
c^Tx
+
\sum_{\omega}
p_\omega
Q_\omega(x).
}
\]

Multistage models must satisfy nonanticipativity:

\[
\boxed{
\text{decisions cannot use future information}.
}
\]

Sample-average approximation replaces an expectation by

\[
\boxed{
\frac1N
\sum_{i=1}^{N}
F(x,\xi^{(i)}).
}
\]

Stochastic gradient descent uses

\[
\boxed{
x_{k+1}
=
x_k
-
\alpha_k
G(x_k,\xi_k).
}
\]

Chance constraints impose reliability:

\[
\boxed{
\mathbb{P}
(
g(x,\xi)\leq0
)
\geq
1-\varepsilon.
}
\]

Risk-averse models can use measures such as

\[
\boxed{
\operatorname{CVaR}_{\alpha}.
}
\]

Stochastic optimization therefore combines

\[
\boxed{
\text{optimization}
+
\text{probability}
+
\text{statistics}
+
\text{simulation}
+
\text{sequential decision making}.
}
\]

%%%%%%%%%%%%%%%%%%%%%%%%%%%%%%%%%%%%%%%%%%%%%%%%%%%%%%%%%%%%
\subsection{Section Connections}
\label{subsec:stochastic-optimization-connections}
%%%%%%%%%%%%%%%%%%%%%%%%%%%%%%%%%%%%%%%%%%%%%%%%%%%%%%%%%%%%

\chapterconnection{
Stochastic Optimization extends the deterministic optimization framework
of Chapter~10 by allowing objectives, constraints, and future decisions to
depend explicitly on uncertainty. Two-stage recourse models connect
directly with linear and convex optimization, while multistage models
connect with Dynamic Programming and stochastic processes from
Chapter~7. Sample-average approximation and stochastic-gradient methods
connect optimization with statistical learning and Monte Carlo methods.
Risk measures such as CVaR link stochastic optimization with finance and
decision theory. The next section, Multiobjective Optimization, develops
problems with several competing objective functions and introduces Pareto
optimality, tradeoff frontiers, scalarization, and compromise solutions.
}

%%%%%%%%%%%%%%%%%%%%%%%%%%%%%%%%%%%%%%%%%%%%%%%%%%%%%%%%%%%%
% SECTION SOURCE TRACKING
%%%%%%%%%%%%%%%%%%%%%%%%%%%%%%%%%%%%%%%%%%%%%%%%%%%%%%%%%%%%

%------------------------------------------------------------
% 10.8 Stochastic Optimization
%------------------------------------------------------------
%
% Source basis:
%   - Standard stochastic programming theory.
%   - Standard stochastic approximation and stochastic-gradient methods.
%   - Standard risk, chance-constrained, and robust optimization theory.
%
% Used for:
%   - stochastic optimization definitions;
%   - uncertain model parameters;
%   - risk-neutral expected objectives;
%   - scenario models;
%   - here-and-now and recourse decisions;
%   - two-stage stochastic programming;
%   - recourse functions;
%   - deterministic equivalents;
%   - nonanticipativity;
%   - multistage stochastic programming;
%   - scenario trees;
%   - stochastic dynamic-programming connections;
%   - wait-and-see solutions;
%   - expected-value solutions;
%   - VSS and EVPI;
%   - complete and relatively complete recourse;
%   - scenario generation and reduction;
%   - sample-average approximation;
%   - law-of-large-numbers connections;
%   - out-of-sample evaluation;
%   - stochastic approximation;
%   - stochastic gradients;
%   - SGD and minibatches;
%   - step-size conditions;
%   - variance reduction preview;
%   - chance constraints;
%   - Gaussian deterministic reformulations;
%   - joint chance constraints;
%   - scenario approximations;
%   - mean--variance optimization;
%   - VaR and CVaR;
%   - coherent risk measures preview;
%   - robust optimization comparison;
%   - distributionally robust optimization;
%   - ambiguity sets;
%   - aleatoric and epistemic uncertainty;
%   - L-shaped decomposition;
%   - SDDP preview;
%   - convexity under expectations;
%   - gradient-expectation interchange;
%   - Jensen's inequality;
%   - decision rules;
%   - rolling-horizon optimization;
%   - probability estimation error;
%   - rare events and importance sampling preview;
%   - applications and exercises.
%
% Notes:
%   - Multiobjective Optimization is developed next in Section 10.9.
%   - Decision Theory is developed in Section 10.14.
%   - Optimal Control is developed in Section 10.16.
%   - Monte Carlo Methods are developed more deeply in Chapter 11.
%
%%%%%%%%%%%%%%%%%%%%%%%%%%%%%%%%%%%%%%%%%%%%%%%%%%%%%%%%%%%%